% Generated draft: requires mathematical and source review.
\documentclass[11pt]{article}
\usepackage[margin=25mm]{geometry}
\usepackage{fontspec}
\setmainfont{Latin Modern Roman}
\usepackage{amsmath,amssymb,mathtools,unicode-math}
\setmathfont{Latin Modern Math}
\usepackage{enumitem,needspace,xurl,hyperref}
\setlength{\parindent}{0pt}
\setlength{\parskip}{5pt}
\setlength{\emergencystretch}{3em}
\newcommand{\N}{\mathbb{N}}
\newcommand{\Z}{\mathbb{Z}}
\newcommand{\Q}{\mathbb{Q}}
\newcommand{\R}{\mathbb{R}}
\newcommand{\C}{\mathbb{C}}
\newcommand{\F}{\mathbb{F}}
\newcommand{\A}{\mathbb{A}}
\newcommand{\PP}{\mathbb{P}}
\newcommand{\E}{\mathbb{E}}
\newcommand{\eps}{\varepsilon}
\newcommand{\dd}{\,\mathrm{d}}
\begin{document}
\section*{UnsolvedMath ID: 701574}
\textbf{Source problem code:} catalogue entry
\subsection{Definitions and mathematical statement}
Let \(G\) be an arbitrary group, and let \(A,B\subseteq G\) be finite subsets with \(|A|\geq 2\) and \(|B|\geq 2\). Define
\[
AB=\{ab:a\in A,\ b\in B\},\qquad
r_{A,B}(x)=\bigl|\{(a,b)\in A\times B:ab=x\}\bigr|
\]
for \(x\in G\), and put
\[
A\mathbin{\cdot_2}B=\{x\in G:r_{A,B}(x)\geq 2\}.
\]
Thus distinct representations mean distinct ordered pairs in \(A\times B\). Define
\[
M(A,B)=\max\left(\{2\}\cup
\bigl\{|gH|:H\leq G,\ g\in G,\ gH\subseteq A\mathbin{\cdot_2}B\bigr\}\right),
\]
where \(gH=\{gh:h\in H\}\) is a left coset. This maximum exists because \(A\mathbin{\cdot_2}B\) is finite, so every qualifying coset is finite and its cardinality is bounded. If no coset qualifies, the maximum is \(2\).

Is it true that, for every such \(G,A,B\),
\[
|A|+|B|-\frac12|AB|-\frac12|A\mathbin{\cdot_2}B|\leq M(A,B)?
\]
The supplied source states that this inequality is proved when \(G\) is abelian. That assertion is retained as a source qualification and has not been independently verified.
\subsection{Short English statement}
For two finite subsets of an arbitrary group, each of size at least two, consider both their product set and the products with at least two distinct representations. The question asks whether the displayed expression in their cardinalities is bounded by the larger of \(2\) and the size of any left coset contained in the latter set. The supplied source reports the abelian case as proved.
\subsection{Sources}
\begin{itemize}
\item Local UnsolvedMath registry, supplied complete record: permanent ID \(701574\), exact source problem code \texttt{catalogue entry}, title \emph{Kourovka Notebook Problem 19.26}. The registry labels the problem \texttt{open}; this is supplied status data, not an independently verified assessment.
\item \emph{The Kourovka Notebook (Unsolved Problems in Group Theory)}, original problem 19.26, 19th issue (2018); the supplied background also identifies the 21st issue edition (2026), \texttt{arXiv:1401.0300v48}. The registry describes the collection as edited by E. I. Khukhro and V. D. Mazurov.
\item The registry identifies W. Dicks as proposer, with proposed year 2018. The supplied statement begins with the separate parenthetical attribution ``(Y. O. Hamidoune).'' Both attributions are retained without assigning an additional role to Hamidoune.
\item The supplied statement reports that the inequality is proved for abelian groups, but supplies no separate proof reference. No external source was visited or independently checked for this draft.
\end{itemize}
\end{document}
